\abstract{%
İteratif yerel aramada (ILS) double-bridge vuruşunun hangi tur kenarlarını kıracağı, gezgin satıcı
probleminde arama verimini belirler. Bu karar için graf sinir ağı (GNN) gibi öğrenilmiş skorlar
önerilmektedir; ancak maliyetlerine değip değmedikleri en basit eğitimsiz alternatife karşı nadiren
sınanır. Hızlı bir 2-opt/Or-opt ILS'te rastgele vuruşu; eğitimsiz, $O(1)$ maliyetli bir kenar
uzunluğu vekilini; LKH-3 turlarıyla eğitilmiş bir GNN'yi; iki çok katmanlı algılayıcıyı (MLP) ve
koşu içinde öğrenen bir modeli karşılaştırdık. Seçimler geliştirme setlerinde yapıldı, protokol
SHA-256 ve zaman damgasıyla donduruldu. Geliştirmede kullanılmamış 54 problemde (10 tohum, 10 sn)
dört ön-kayıtlı hipotezin dördü de desteklendi: vekil, rastgele vuruşa göre gap'i 0,172 puan
düşürdü ($p=3{,}7\times10^{-7}$) ve GNN ile kazanç hedefli MLP'ye karşı 0,05 puanlık sınırla
kötü-değil testlerini geçti. Önceden kaydedilen keşifsel bloklarda vekil, ağın süresi bütçeden
düşüldüğünde yayımlanmış NeuroLKH modelinden ayrışmadı; süre sayılmadığında NeuroLKH problemlerin
çoğunda küçük bir farkla öndeydi. Eşit vuruş sayısında öğrenilmiş skorlar vekilden kötü değildi ve
vuruş başına en iyi sinyal eğitimsiz kNN sırasıydı: eşit süredeki sıralamayı sinyalin niteliğinden
çok maliyeti belirler. Çevrimiçi modelin zararı işaretli ödülden geliyordu; kabule dayalı krediyle
kayboldu. LKH-3 içinde hedeflemenin ölçülebilir bir etkisi olmadı. Hızlı ILS'te vuruş hedefleme
öğrenilmiş model gerektirmedi; öğrenilmiş arama bileşenleri önce en basit eğitimsiz vekillere
karşı ve eşit maliyetle sınanmalıdır.
}
\keywords{gezgin satıcı problemi, iteratif yerel arama, vuruş hedefleme, eğitimsiz vekil, ön-kayıtlı deney, NeuroLKH}
% kaynak: E584, E585, E589, E599-E602, E1106, E1107, E1128-E1131, E1148-E1152, E1172
